% concatenated from main.tex
%\documentclass[a4paper,10pt]{article}
\documentclass[a4paper]{article}

\usepackage[colorlinks=true, allcolors=blue]{hyperref}
%\usepackage{epstopdf} %to include eps images
%\epstopdfsetup{update}
\usepackage{authblk}
\usepackage{bbm}
%\usepackage{comment}
\usepackage{overpic}
\usepackage{rotating}
%\usepackage[numbers,sort]{natbib}
\usepackage{xcolor}
 \usepackage{relsize}


% standard packages
\usepackage[utf8]{inputenc}
\usepackage[english]{babel}
\usepackage{csquotes}
\usepackage{geometry}
\usepackage{amsfonts}
\usepackage{amsmath}
\usepackage{amsthm, thmtools}
\usepackage{amssymb}
\usepackage{xcolor}
\usepackage{bm}



\usepackage{algpseudocode}
\usepackage{algorithm}
\usepackage{algorithmicx}
\usepackage{amsopn}



%\usepackage[pdftex]{graphicx}
\usepackage{caption}
\usepackage{subcaption}

\DeclareGraphicsExtensions{.png,.pdf,.jpg}

\DeclareMathOperator{\Id}{Id}


\def\rmd{{\mathrm{d}}}

\numberwithin{equation}{section}

\newcommand{\gom}[1]{\textcolor{purple}{GOM: #1}}
\newcommand{\me}[1]{\textcolor{blue}{ME: #1}}

\newcommand{\TU}[1]{\textup{#1}}

%\newtheorem{defn}[section]{Definition}
%\newtheorem{prop}[section]{Proposition}


\newtheorem{axiom}{Axiom}
\newtheorem{claim}[axiom]{Claim}
\newtheorem{theorem}{Theorem}[section]
\newtheorem{lemma}[theorem]{Lemma}
\newtheorem{proposition}[theorem]{Proposition}
\newtheorem{corollary}[theorem]{Corollary}
\newtheorem{conjecture}[theorem]{Conjecture}
\newtheorem{definition}[theorem]{Definition}
\newtheorem{remark}[theorem]{Remark}
\newtheorem{notations}[theorem]{Notations}
\newtheorem{example}[theorem]{Example}


\newcommand{\bibfilename}{mrabbrev,mybibfile} % <- replace 'mybibfile' with basename of your bibfile

\newtheorem{goal}{Goal}

\begin{document}
\title{Affine Invariant Langevin Dynamics for rare-event sampling}

%\author[1]{Deepyaman Chakraborty\thanks{deepyaman.chakraborty@fu-berlin.de}}
\author[1]{Deepyaman Chakraborty}

%\author[1]{Ruben Harris\thanks{r.harris@fu-berlin.de}}
\author[1]{Ruben Harris}

%\author[1]{Rupert Klein\thanks{rupert.klein@math.fu-berlin.de}}
\author[1]{Rupert Klein}

\author[2]{Guillermo Olic\'on-M\'endez}
%\author[2]{Guillermo Olic\'on-M\'endez\thanks{guillermo.olicon@uni-potsdam.de}}

%\author[2]{Sebastian Reich\thanks{sereich@uni-potsdam.de}}
\author[2]{Sebastian Reich}


%\author[1]{Claudia Schillings\thanks{c.schillings@fu-berlin.de}}
\author[1]{Claudia Schillings}


\affil[1]{Freie Universit\"at Berlin, Institut f\"ur Mathematik und Informatik, 14195 Berlin, Germany}

\affil[2]{Universit\"at Potsdam, Institut f\"ur Mathematik, 14476 Potsdam, Germany}


\maketitle


\begin{abstract}
We introduce an affine invariant Langevin dynamics (ALDI) framework for the
efficient estimation of rare events in nonlinear dynamical systems. Rare events are formulated as Bayesian inverse
problems through a nonsmooth limit-state function whose zero level set characterises the
event of interest. To overcome the nondifferentiability of this function, we
propose a smooth approximation that preserves the failure set and yields a
posterior distribution satisfying the small-noise limit. The resulting potential
is sampled by ALDI, a (derivative-free) interacting particle system whose affine
invariance allows it to adapt to the local anisotropy of the posterior.

We demonstrate the performance of the method across a hierarchy of benchmarks, namely
two low-dimensional examples (an algebraic problem with convex geometry and a dynamical problem of saddle-type instability) and a point-vortex model for atmospheric blockings. In all cases, ALDI
concentrates near the relevant near-critical sets and provides accurate
proposal distributions for self-normalised importance sampling. The framework is
computationally robust, potentially gradient-free, and well-suited for complex forward
models with strong geometric anisotropy. These results highlight ALDI as a
promising tool for rare-event estimation in unstable regimes of dynamical systems.
\end{abstract}


% \begin{keywords}
%   Bayesian inverse problems, rare-event estimation, Affine invariant Langevin dynamics, point-vortex dynamics
% \end{keywords}


% \begin{AMS}
% 65C05, 65C30, 62F15, 76B47
% \end{AMS}



{\small \textbf{Keywords:}  Bayesian inverse problems, rare-event estimation, Affine invariant Langevin dynamics, point-vortex dynamics}

{\small \textbf{2020 Mathematics subject classification:} 65C05, 65C30, 65C35, 62F15, 76B47}



\section{Introduction}
\label{sec:intro}
Rare events play a crucial role in many complex dynamical systems, ranging from transitions between metastable states in molecular dynamics \cite{KoltaiEtAl2016} to critical shifts in atmospheric and oceanic circulation regimes. In geophysical contexts, such events include persistent \emph{atmospheric blockings}---large-scale quasi-stationary flow patterns that interrupt the usual westerly circulation and can lead to prolonged periods of heat or precipitation extremes. Although these phenomena occur infrequently, their consequences can be severe, and quantifying their probability and underlying mechanisms remains a major scientific challenge \cite{DaviniEtAl2021,Hirt18,lucarini2019,Muller15,Ragone20}

The mathematical and computational treatment of rare events requires specialised tools, since direct numerical simulation becomes infeasible due to the exponentially small probability of observing such anomalies. Classical approaches include importance sampling and splitting methods \cite{Beck2017, cerou}, variational methods \cite{hartmann}, large deviation theory \cite{stadler} and particle methods \cite{delmoral}. More recently, a probabilistic reformulation of rare-event estimation as a \emph{Bayesian inverse problem} has been proposed \cite{Wagner21}, enabling the use of sampling and inference techniques from uncertainty quantification. In this perspective, the rare event is identified with the conditioning of a prior measure on the event $\{G(u)\le 0\}$, where $G$ is a suitably chosen observable. The resulting posterior is highly concentrated in a low-probability region, making it a challenging target for standard sampling algorithms.

Ensemble-based methods such as the Ensemble Kalman Filter (EnKF) have been widely used in data assimilation and inverse problems due to their computational efficiency and ability to handle high-dimensional state spaces \cite{evensen2009,reichcotter2015}. They propagate and update an ensemble of particles according to system dynamics and observational constraints, providing a derivative-free approximation of Bayesian updates. Ensemble-based methods such as the EnKF have proved effective far beyond the linear–Gaussian setting and are routinely applied to nonlinear, non-Gaussian problems \cite{evensen2003,evensen2009,reichcotter2015,vanleeuwen2010}. Their analysis step, however, is built around updating the first and second moments via a linear correction, which—together with finite-ensemble effects, localisation, and inflation—can bias tail probabilities and collapse separated modes unless special care is taken (e.g., tempering, iterative updates, or mixture/particle–ensemble hybrids) \cite{vanleeuwen2010,bocquet2010}. Consequently, representing an extremely small posterior mass in the tail region relevant for rare events remains challenging in practice, even though many EnKF variants alleviate parts of this issue.

To target rare events more directly, \cite{Wagner21} proposed an EnKF-based formulation in which the event is encoded through an observable $G$ (taken, without loss of generality, to be nonpositive on the event), and the inverse problem
\begin{equation}
  \tilde{G}(u)=0, 
  \qquad \text{with} \quad
  \tilde{G}(u):=\max\{0,G(u)\},
\end{equation}
so that the event corresponds to the set $\{G(u)\le 0\}$. This casts rare-event estimation as Bayesian conditioning on (a smoothed version of) the indicator of the rare set, thereby providing an inference route that we will retain here while replacing the ensemble update with an affine-invariant Langevin sampler (ALDI) \cite{Nusken20} as a stochastic, geometry-adapted sampling method for rare-event estimation. ALDI combines affine invariance, ensuring that the sampling dynamics is unaffected by linear rescaling or anisotropy of the state space, with the stochastic exploration of Langevin-type diffusions. This results in improved stability and efficiency, especially in high-dimensional or strongly anisotropic settings. The present work thus bridges ensemble-based inference, stochastic dynamics, and rare-event analysis within a unified geometric framework, building upon earlier developments in Bayesian sampling and dynamical systems theory. In \cite{beh2025affineinvariantinteractinglangevin}, a related method based on a combination of Langevin samplers and importance sampling (IS) has been introduced, which employs a smooth
approximation of the failure event and provides an analysis of both the resulting
importance sampling error and the discretization bias of the Langevin algorithm used for sampling the smoothed density. This forms the basis for an optimal computational work distribution between the Langevin sampler and the IS step. In contrast, our focus is on the small-noise limit of the
associated Bayesian formulation: we characterise the limiting posterior as the prior
restricted to the failure set and prove total variation convergence. This offers a
measure-theoretic foundation that is independent of any specific sampling scheme and
complements the algorithmic analysis in~\cite{beh2025affineinvariantinteractinglangevin}.

We apply ALDI to rare-event estimation in two simple test cases to analyse its behaviour, and then a point-vortex model for two-dimensional flow, used to explore the persistence of quasistationary vortex configurations associated with atmospheric blockings \cite{Kuhlbrodt00}.
In both cases, we identify two qualitatively different mechanisms that generate low-probability events:
\begin{itemize}
    \item \textit{Parametric rare events}, where the event $\{G_\lambda(u)\le 0\}$ occurs only for specific parameter values $\lambda$, and
    \item \textit{Dynamic rare events}, where the event arises from the intrinsic instability of the system’s dynamics near invariant structures.
\end{itemize}

Understanding and quantifying such rare transitions is essential in many applications, as they often correspond to critical transitions or bifurcations in stochastic systems. Rare events are also related to bifurcations in stochastic dynamics when the noisy perturbations are bounded \cite{olicon24,olicon25}. Atmospheric blockings, for instance, represent prolonged excursions of the flow dynamics toward atypical attractors, the occurrence of which is often controlled by a delicate interplay between stochastic forcing, energy input, and multiscale interactions \cite{klein2021,ragone2021}.

\paragraph{Contributions and structure}
The main contributions of this paper are as follows:
\begin{itemize}
    \item Based on the reformulation of rare-event estimation as a Bayesian inverse problem, we introduce an ALDI method for exploring the associated conditional measure.
    \item We prove that, as the observation noise tends to zero, the estimator obtained from the affine-invariant Langevin dynamics converges to the true probability of the rare event, independently of the value of the smoothing parameter.
    \item We demonstrate the performance of ALDI on dynamical systems around unstable invariant objects, including vortex configurations related to atmospheric blockings.
\end{itemize}
The remainder of this paper is structured as follows.
Section~\ref{SEC:preliminaries} reviews the Bayesian formulation of rare-event estimation and presents the affine-invariant Langevin dynamics. To exploit its theoretical features, we introduce a smoothing technique used in this context, along theoretical results concerning the invariance and convergence of the method. In Section~\ref{SEC:numerics}, we test the method on two simple problems: an algebraic convex example (Subsection~\ref{SUBSEC:convex}) and a linear unstable dynamical system (Subsection~\ref{SUBSEC:saddle}). We further exploit the method for rare-event estimation in point-vortex dynamics (Subsection~\ref{SUBSEC:vortex}). Finally, Section~\ref{SEC:conclusions} summarises the results and discusses perspectives for future research.



\section{Preliminaries}
\label{SEC:preliminaries}
In this paper, we consider rare events defined by the outcome of a \textit{limit-state function} $G$. More precisely, let $(\Omega,\mathcal{F},\mathbb{P})$ be a probability space, and the state of a system given by a random variable $X:\Omega\rightarrow \mathbb{X}$, where the state space $\mathbb{X}=\mathbb R^d$ endowed with its Borel $\sigma$-algebra $\mathcal{B}$, and distributed according to $\mu_0(\mathrm d x)=\rho_0(x) \mathrm dx$. 

We will be interested in the convergence of a family of measures $\mu_\delta$ on $(\mathbb X, \mathcal{B})$ to a given measure $\mu$. Recall that the \emph{total variation distance} between probability distributions $\mu_1,\mu_2$ is defined as $\Vert \mu_1-\mu_2\Vert_{TV}:= \sup_{A\in\mathcal{B}} \vert \mu_1(A)-\mu_2(A)\vert$. Moreover, when both distributions are absolutely continuous with respect to the same distribution $\mu_0$, so that $\mu_i(dx)=\rho_i(x)\mu_0(dx)$, then
\[\Vert \mu_1-\mu_2\Vert_{TV}\ = \frac{1}{2}\Vert \rho_1-\rho_2\Vert_{L^1(\mu_0)}:= \frac{1}{2}\int_{\mathbb X} \vert \rho_1-\rho_2\vert d\mu_0. \]
Given a density $\rho$ of a distribution $\mu$, we will sometimes refer to the Lebesgue space $L^1(\rho)\equiv L^1(\mu)$.


\subsection{Rare events as Bayesian inverse problems}

For a given \emph{limit-state function} \( G : \mathbb{X} \rightarrow \mathbb{R} \), the \emph{failure domain} is defined (without loss of generality) by
\[
F=\{ G(x) \le 0 \},
\]
and we are interested in estimating the associated probability of failure
\begin{equation}
    \label{eq:probfailure_event}
    P_f := \mathbb{P}\left( \left\{ \omega \in \Omega : G(X(\omega)) \le 0 \right\} \right)=\int_\mathbb{X} \mathbf{1}_{F}(x) \mu_0(\mathrm dx)
\end{equation}
where $\mathbf 1 $ denotes the indicator function and $\mu_0$ a given measure on $\mathbb X$. 
In many situations, the failure domain is highly nonlinear or disconnected, and the corresponding event is rare. Consequently, estimating \( P_f \) accurately requires specialised methods that go beyond standard Monte Carlo techniques, which tend to undersample low-probability regions. One fruitful approach is to recast the identification of failure states as an \emph{inverse problem}. 

To this end, following \cite{Wagner21}, we introduce the modified limit-state function
\[
\tilde{G}(x) := \max\{ 0, G(x) \},
\]
and consider the inverse problem
\begin{equation}
    \label{eq:inverse}
    \tilde{G}(x) = 0,
\end{equation}
whose solution set coincides with the failure domain $\{G(x)\le 0\}$.

We adopt a Bayesian formulation by introducing an artificial observation model with additive noise,
\begin{equation}
    \label{eq:bayesian_inverseproblem}
    y = \tilde{G}(x) + \eta,
\end{equation}
where \( \eta \sim \mathcal{N}(0, R) \) with variance parameter \( R > 0 \), and we fix the observed value to \( y = 0 \). The noise variance \( R \) acts as a \emph{smoothing parameter}: as \( R \to 0 \), the posterior concentrates near the zero level set of \( \tilde{G} \), i.e., the boundary of the failure region.

Denoting by \( \rho_0(x) \) the density of $\mu_0$, Bayes’ theorem yields the posterior density
\begin{equation}
    \label{eq:posterior distribution}
    \rho_*(x) = \frac{1}{Z} \exp\!\left(-\frac{1}{2R}\tilde{G}(x)^2\right)\rho_0(x),
\end{equation}
where \( Z \) is the normalising constant. The posterior distribution is in general non-Gaussian, since \( \tilde{G} \) is typically nonlinear and may only be continuous. In the following we assume that the prior density $\rho_0$ is strictly positive
and continuous on $\mathbb X$, and that the limit-state function $G:\mathbb X\to\mathbb R$ is
continuous. Consequently, $\tilde G(x)=\max\{0,G(x)\}$ is also continuous but
not differentiable on the boundary $\{G(x)=0\}$, which motivates the
smoothing procedure described in Section~\ref{subsubsec:smooth}.

In \cite{Wagner21}, a variant of the EnKF was employed to approximate the failure probability \(P_f\) based on this Bayesian reformulation. In the present work, we replace the EnKF by ALDI, whose invariant measure coincides with the smoothed posterior \eqref{eq:posterior distribution}. Unlike the EnKF, which relies on affine ensemble updates and provides a Gaussian approximation, ALDI preserves the full non-Gaussian structure of the posterior and remains consistent in the small-noise and vanishing-smoothing limits.

\subsection{Affine-Invariant Interacting Langevin Dynamics}\label{subsec:ALDI}

In our numerical experiments, we employ ALDI \cite{Nusken20} to sample the posterior distribution associated with the Bayesian inverse problem \eqref{eq:bayesian_inverseproblem}. In general, ALDI samples from a probability density of the form
\[
\rho_*(x) = Z^{-1}\exp(-\Phi(x)),
\]
where \(Z>0\) is a normalisation constant and \(\Phi:\mathbb{R}^d\to\mathbb{R}\) denotes the \emph{potential function}.  
In our setting, the potential function reads
\begin{equation}
    \label{eq:potential_function}
    \Phi(x) = \frac{1}{2R}\,\tilde{G}(x)^2 - \ln \rho_0(x),
\end{equation}
so that the distribution corresponding to $\rho_*$ is absolutely continuous with respect to the prior measure \(\mu_0\), with density proportional to \(\exp(-\tfrac{1}{2R}\tilde{G}(x)^2)\). We assume $\rho_0$ to be strictly positive and smooth enough so that $-\ln \rho_0$ is well-defined and differentiable where needed.

The ALDI method is formulated as an interacting particle system.  
If the potential function \(\Phi\) is smooth, the dynamics of an ensemble of \(J\) particles \(X=(x^{(1)},\ldots,x^{(J)})\) is given by the system of Itô-type stochastic differential equations
\begin{equation}
    \label{eq:ALDI_smooth}
    dx_t^{(j)} 
    = \left[-\mathcal{C}(X_t)\,\nabla_{x^{(j)}} \Phi\!\left(x_t^{(j)}\right)
    + \frac{d+1}{J}\big(x_t^{(j)} - m(X_t)\big)\right] dt
    + \sqrt{2}\,\mathcal{C}^{1/2}(X_t)\,dW_t^{(j)},
\end{equation}
where \(W_t^{(j)}\) are independent \(d\)-dimensional standard Brownian motions.  
Here, \(m(X)\) and \(\mathcal{C}(X)\) denote the empirical mean and covariance of the ensemble,
\begin{equation}
    \label{eq:ALDI_definitions}
    \begin{split}
        m(X) &:= \frac{1}{J}\sum_{j=1}^J x^{(j)}, \\[2mm]
        \mathcal{C}(X) &:= \frac{1}{J}\sum_{j=1}^J \big(x^{(j)}-m(X)\big)\big(x^{(j)}-m(X)\big)^\top, \\[2mm]
        \mathcal{C}^{1/2}(X) &:= \frac{1}{\sqrt{J}}\big(X - m(X)1_J\big),
    \end{split}
\end{equation}
with \(1_J=(1,\ldots,1)^\top \in \mathbb{R}^J\).

Let \(\pi_t^{(J)}\) denote the probability density of \(X_t\).  
The corresponding extended invariant density is
\begin{equation}
    \label{eq:invariant_density}
    \pi_*^{(J)}(x) = \prod_{j=1}^J \rho_*\!\left(x^{(j)}\right).
\end{equation}

\begin{theorem}
\label{THM:ALDI_properties}
Let $\mathbb X=\mathbb{R}^d$. Assume that the initial covariance \(\mathcal{C}(X_0)\) is positive definite, the potential \(\Phi \in C^2(\mathbb{R}^d)\cap L^1(\rho_*)\), and that there exist a compact set \(K \subset \mathbb{R}^d\) and constants \(0 < c_1 < c_2\) such that for all \(x \in \mathbb{R}^d \setminus K\),
\begin{equation}
    \label{eq:conditions_potential}
    \begin{split}
        c_1 \|x\|^2 \le \Phi(x) \le c_2 \|x\|^2, \quad
        c_1 \|x\| \le \|\nabla \Phi(x)\| \le c_2 \|x\|, \quad
        c_1 I_d \le \mathrm{Hess}\,\Phi(x) \le c_2 I_d.
    \end{split}
\end{equation}
Then the ALDI system \eqref{eq:ALDI_smooth} satisfies the following properties:

\smallskip
\noindent
1. \textbf{(Affine invariance)}  
The nonlinear Fokker–Planck equation associated with \eqref{eq:ALDI_smooth},
\begin{equation}
\label{eq:Fokker_Planck_nonlinear}
\partial_t\pi_t = 
\nabla_x \!\cdot\!
\left(
\pi_t\,C(\pi_t)\,\nabla_x
\frac{\partial \mathrm{KL}(\pi_t\Vert \rho_*)}{\partial \pi_t}
\right),
\quad
C(\pi_t) = \mathbb{E}_{\pi_t}\!\left[(x-\mu_t)(x-\mu_t)^\top\right],
\end{equation}
is affine invariant.

\smallskip
\noindent
2. \textbf{(Invariant distribution)}  
The product measure \(d\mu(x)=\pi_*^{(J)}(x)\,dx\) is invariant for the particle system \eqref{eq:ALDI_smooth}.

\smallskip
\noindent
3. \textbf{(Global strong solution)}  
The system \eqref{eq:ALDI_smooth} admits a unique global strong solution \(X_t\), and the empirical covariance \(\mathcal{C}(X_t)\) remains strictly positive definite almost surely.

\smallskip
\noindent
4. \textbf{(Ergodicity)}  
If \(J > d+1\), then \(\pi_t^{(J)} \to \pi_*^{(J)}\) in total variation distance as \(t \to \infty\).

\end{theorem}

\begin{proof}
We refer to Lemma~2.4, Corollary~4.2, Proposition~4.4, Proposition~4.5, and Lemma~4.7 in \cite{Nusken20}.
\end{proof}

\subsubsection{A smoothing of the limit-state function}\label{subsubsec:smooth}
As mentioned earlier, our problem involves a potential $\Phi$ that is nonsmooth. In order to exploit the theoretical aspects of the ALDI method \eqref{eq:ALDI_smooth} as given in Theorem~\ref{THM:ALDI_properties}, we introduce a smooth approximation of $\tilde{G}$ which coincides with it for any value of $G(x)$ outside a small interval $(0,\delta)$. Consider the function
\begin{equation}
\label{eq:psi_base}
\psi_\delta(x):=\begin{cases}
\exp\left(\frac{1}{\delta^2}\right)\cdot \exp\left(-\frac{1}{x^2}\right), & \textup{if }
x> 0\\
0, & \textup{otherwise}, \end{cases}
\end{equation}
and the ramp function
\begin{equation}
    \label{eq:ramp_function}
\phi_\delta(x):=\frac{\psi_\delta(x)}{\psi_\delta(x)+ \psi_\delta(\delta-x)},
\end{equation}
which approximates smoothly the indicator function $\mathbf{1}_{[0,\infty)}$. Therefore, a smooth approximation of $\tilde{G}$ is given by
\begin{equation}
    \label{eq:smooth_G}
\tilde G_\delta(x) :=
\begin{cases}
  0, & \text{if } G(x) < 0,\\[0.3em]
  G(x)\dfrac{\psi_\delta(G(x))}{\psi_\delta(G(x))+\psi_\delta(\delta-G(x))}, & \text{if } G(x)\in[0,\delta],\\[0.7em]
  G(x), & \text{if } G(x) > \delta.
\end{cases}
\end{equation}
Note that $\tilde G_\delta(x)\to\tilde G(x)$ uniformly as $\delta\to0$, and the zero level set $\{\tilde{G}=0\}$ is preserved for all $\delta\geq0$. 


\subsubsection{Consistency of the smoothed Posterior and ALDI}
\label{SUBSUBSEC:consistency}
We now establish that the smoothed posterior distribution underlying the ALDI dynamics 
converges, in the joint limit of vanishing observation noise and vanishing smoothing, 
to the prior distribution restricted to the rare-event domain.

\begin{theorem}[Consistency of the smoothed posterior]
\label{thm:consistency}
Let $F=\{x\in\mathbb X : G(x)\le 0\}$ be the failure domain and assume that $P_f=\mu_0(F)>0$. Let $\tilde G_\delta$ be the smoothed approximation defined in
\eqref{eq:smooth_G} and, for each $R>0$, define the probability measures
\begin{equation}\label{eq:smoothedpost}
  \mu_{\delta,R}(dx)
  := \frac{1}{Z_{\delta,R}}
     \exp\!\left(-\frac{\tilde G_\delta(x)^2}{2R}\right)\rho_0(x)\,dx,
  \quad  
  Z_{\delta,R}
  := \int_{\mathbb X} \exp\!\left(-\frac{\tilde G_\delta(x)^2}{2R}\right)
      \rho_0(x)\,dx.
\end{equation}
Then, for every fixed $\delta>0$ we have $\mu_{\delta,R}\;\xrightarrow{\mathrm{TV}}\;\mu_F$ as $R\rightarrow 0 $,
where the distribution $\mu_F(dx)=\frac{\mathbf{1}_F(x)}{\mu_0(F)}\rho_0(x)dx$ is independent of~$\delta$. 
\end{theorem}

\begin{proof}

Consider $\delta>0$ fixed. First, we show that the normalisation constants $Z_{\delta,R}$ converge to $\mu_0(F)$. Indeed, recall from \eqref{eq:smooth_G} that $\tilde{G}_\delta(x)=0$ whenever $x\in F$, for all $\delta\geq 0$. By splitting the state space $\mathbb{X}=F \cup (\mathbb X \setminus F)$, it follows from \eqref{eq:smoothedpost} that
\[
    Z_{\delta,R} = \mu_0(F) + \int_{\mathbb X \setminus F} \exp\left( -\frac{\tilde{G}_\delta(x^2)}{2R} \right) \rho_0(x)dx.
\]
The integrand on the right hand side is bounded by $\rho_0(x)$, and $\exp(-\tilde{G}(x)^2/(2R))\rightarrow 0$ for each $x\in \mathbb{X}\setminus F$. Hence, by the dominated convergence theorem we obtain that
\begin{equation}
    \label{eq:limit_normalisation}
        \lim_{R\rightarrow 0} Z_{\delta,R}=\mu_0(F).
\end{equation}

Denote $f_{\delta,R}(x):=\frac{\exp\left(- \tilde{G}(x)^2/Z_{\delta,R}\right)}{2R}\equiv \frac{N_{\delta,R}(x)}{Z_{\delta,R}}$ the density of $\mu_{\delta,R}$ with respect to $\mu_0$. We show that $f_{\delta,R}\rightarrow f_F(x):= \frac{1}{\mu_0(F)}\mathbf{1}_F(x)$ in $L^1(\mu_0)$. Indeed, by splitting again the state space and recalling that $\tilde G(x)=0$ for $x\in F$, we get
\begin{align*}
    \int_{\mathbb X}\vert f_{\delta,R}-f_F \vert d\mu_0 & = \int_F \left\vert \frac{N_{\delta,R}(x)}{Z_{\delta,R}} - \frac{1}{\mu_0(F)}\right\vert d\mu_0(x) + \int_{\mathbb X \setminus F} \frac{N_{\delta, R}(x)}{Z_{\delta,R}}d\mu_0(x)= \\
    &= \frac{\left\vert Z_{\delta,R}-\mu_0(F) \right\vert}{Z_{\delta,R}}+ \frac{1}{Z_{\delta,R}}\int_{\mathbb{X}\setminus F}N_{\delta,R}(x)d\mu_0(x).
\end{align*}
As $R\rightarrow 0$, the first term on the right hand side vanishes as given in \eqref{eq:limit_normalisation}. For the second term, observe that $N_{\delta,R}\rightarrow \mathbf{1}_F$ pointwise and that $N_{\delta,R}\leq 1 \in L^1(\mu_0)$, so that the integral on the right hand side vanishes as $R\rightarrow 0$ due to the dominated convergence theorem. 

Since the densities $f_{\delta,R}$ converge to $f_F$ in $L^1(\mu_0)$, their corresponding measures $\mu_{\delta,R}$ converge to $\mu_F$ in total variation as they are all absolutely continuous with respect to $\mu_0$, and the result follows.
\end{proof}



Theorem~\ref{THM:ALDI_properties} can be translated in terms of the ALDI method as given in the following statement.

\begin{corollary}[Consistency of ALDI]
For every fixed smoothing parameter $\delta>0$ and every noise level $R>0$,
the ALDI dynamics \eqref{eq:ALDI_smooth} admits the invariant product measure
\[
\pi^{(J),\delta,R}_{*}(x) \;\propto\; 
\prod_{j=1}^J 
\exp\!\left(-\frac{ G_\delta(x^{(j)})^2 }{2R}\right)\rho_0(x^{(j)}),
\]
and is ergodic whenever $J>d+1$ (see Theorem~\ref{THM:ALDI_properties}).  
If the dynamics is run to stationarity, each particle satisfies
\[
X^{(j)}_{\delta,R}\sim \rho^*_{\delta,R},
\]
where $\rho^*_{\delta,R}$ is the smoothed posterior defined in
\eqref{eq:smoothedpost}.

Moreover, by Theorem~\ref{thm:consistency}, for every fixed $\delta>0$,
\[
\rho^*_{\delta,R} 
\;\xrightarrow[R\to 0]{\rm TV}\;
\mu_F ,
\qquad 
F=\{x : G(x)\le 0\},
\]
and the limit measure $\mu_F$ is independent of $\delta$.
Consequently, the stationary ALDI estimator is consistent for the
rare-event distribution: for each ensemble index $j$,
\[
X^{(j)}_{\delta,R}
\;\xrightarrow[R\to 0]{\rm law}\;
\mu_F .
\]
\end{corollary}


\begin{remark}
The proof of Theorem~~\ref{thm:consistency} relies only on the continuity of $G$, 
the property that $\tilde G_\delta(x)=0$ if and only if $G(x)\le 0$.
In particular, for every fixed $\delta>0$ we have
\[
\mu_{\delta,R} \xrightarrow[R\to 0]{\mathrm{TV}} \mu_F,
\]
and the limit measure $\mu_F$ is independent of $\delta$. 
Thus, $\delta$ acts purely as a smoothing parameter that improves the 
regularity of the potential but does not affect the limiting conditional
distribution as $R\to 0$. 

The result itself is formulated at the level of the target measures 
$\mu_{\delta,R}$ and is therefore independent of the particular sampling 
scheme. In particular, any algorithm whose invariant distribution is 
$\mu_{\delta,R}$ (such as the gradient-based ALDI dynamics) inherits the same consistency property.
\end{remark}

 \subsection{Importance sampling for the probability of failure}
\label{subsec:IS}

Recall that our goal is to estimate the probability of failure
\[
  P_f = \mu_0(F)
      = \int_X \mathbf{1}_{\{G(x)\le 0\}}\,\rho_0(x)\,dx,
  \qquad F := \{x\in X : G(x)\le 0\}.
\]
A standard approach is IS, which rewrites
\[
  P_f
  = \int_F \frac{\rho_0(x)}{q(x)}\, q(x)\,dx,
\]
where $q$ is a proposal density. The usual global support condition
$\operatorname{supp}(\rho_0)\subseteq \operatorname{supp}(q)$ is not required here.
Since the integrand contains $\mathbf{1}_{F}$, the correct condition is only
\begin{equation}
  \label{eq:support-cond-F}
  \operatorname{supp}\bigl(\mathbf{1}_{F}\rho_0\bigr)
    \subseteq \operatorname{supp}(q),
\end{equation}
i.e., $q(x)>0$ for all $x\in F$ with $\rho_0(x)>0$.
This is naturally satisfied for the constructions below. A central difficulty in importance sampling is \emph{weight degeneracy}, where only a small number of samples carry non-negligible weights due to a mismatch between the proposal distribution and the rare-event region, resulting in large estimator variance and a severely reduced effective sample size. In particular, importance sampling based on the prior as proposal is highly prone to weight degeneracy in rare-event regimes, whereas ALDI-based proposals concentrate mass near the failure set and therefore yield substantially more balanced weights.

We consider two IS strategies based on stationary samples from
$\rho^*_{\delta,R}$, an auxiliary smoothed density used to bias sampling
towards the failure set.
Both methods are consistent for estimating the failure probability $P_f$,
but differ in how the proposal distribution is constructed and exploited.

\paragraph{Variant~1: Product estimator based on ALDI samples.}
Assume that the samples $\{X_m\}_{m=1}^M$ are obtained from ALDI at stationarity,
so that $X_m\sim \rho^*_{\delta,R}$, where
\[
  \rho^*_{\delta,R}(x)
  = \frac{1}{Z^*_{\delta,R}}\,
    \exp\!\Bigl(-\tfrac{1}{2R}\tilde G_\delta(x)^2\Bigr)\,\rho_0(x),
\]
and
\[
  Z^*_{\delta,R}
  := \int \exp\!\Bigl(-\tfrac{1}{2R}\tilde G_\delta(x)^2\Bigr)\,\rho_0(x)\,dx .
\]

Since $\tilde G_\delta(x)=0$ for all $x\in F=\{G(x)\le 0\}$, the failure
probability can be written as
\[
P_f
= \int_F \rho_0(x)\,dx
= Z^*_{\delta,R}\,
  \mathbb P_{\rho^*_{\delta,R}}(X\in F).
\]
This identity motivates the product estimator
\[
  \widehat P_f
  = \widehat Z^*_{\delta,R}\,
    \widehat p^*,
  \qquad
  \widehat p^*
  := \frac{1}{M}\sum_{m=1}^M
      \mathbf{1}_{\{G(X_m)\le 0\}},
\]
where $\widehat p^*$ is the empirical probability of the failure event under
$\rho^*_{\delta,R}$, and $\widehat Z^*_{\delta,R}$ is an estimator of the
normalisation constant. Since $\rho^*_{\delta,R}(x)=\frac{1}{Z^*_{\delta,R}}\exp\!\big(-\tfrac{1}{2R}\tilde G_\delta(x)^2\big)\,\rho_0(x)$, we can rewrite
\[
\rho_0(x)=Z^*_{\delta,R}\,\exp\!\big(\tfrac{1}{2R}\tilde G_\delta(x)^2\big)\,\rho^*_{\delta,R}(x).
\]
Integrating both sides over $x$ yields the inverse-moment identity
\[
1 = Z^*_{\delta,R}\,\mathbb E_{\rho^*_{\delta,R}}\!\left[\exp\!\big(\tfrac{1}{2R}\tilde G_\delta(X)^2\big)\right],
\qquad\Rightarrow\qquad
Z^*_{\delta,R}=\left(\mathbb E_{\rho^*_{\delta,R}}\!\left[\exp\!\big(\tfrac{1}{2R}\tilde G_\delta(X)^2\big)\right]\right)^{-1},
\]
which motivates the Monte--Carlo approximation in Alg.~2.1 by replacing the expectation with an empirical average over stationary ALDI samples.

\begin{algorithm}[t]
\caption{ALDI-based product estimator for $P_f$}
\label{alg:aldi_product}
\begin{algorithmic}[1]
\Require prior $\rho_0$, limit-state $G$, smoothing parameters $\delta,R$,
         ALDI integrator, ensemble size $J$, sample size $M$
\Ensure estimate $\widehat P_f$
\State initialise ALDI ensemble
\State run ALDI (possibly with annealing) towards $\rho^*_{\delta,R}$
\State collect $M$ approximately independent samples $\{X_m\}$
\State compute
       $\displaystyle
        \widehat p^*=
        \frac{1}{M}\sum_{m=1}^M \mathbf{1}_{\{G(X_m)\le 0\}}$
\State estimate normalisation constant $\widehat Z^*_{\delta,R}
=
\left(
\frac{1}{M}\sum_{m=1}^M
\exp\!\Bigl(\tfrac{1}{2R}\tilde G_\delta(X_m)^2\Bigr)
\right)^{-1}.
$
       
\State compute $\widehat P_f=\widehat Z^*_{\delta,R}\,\widehat p^*$
\State \Return $\widehat P_f$
\end{algorithmic}
\end{algorithm}

\begin{remark}
Alg.~\ref{alg:aldi_product} is included for conceptual completeness. We will see in the numerical experiments that the normalisation estimator is numerically unstable, as given in Figure~\ref{fig:size_and_failed}(a).
\end{remark}

\paragraph{Variant~2: Mixture-based IS using a fitted proposal.}
To increase the accuracy of the estimator to a prescribed tolerance, one
typically needs to generate a large number of independent samples from the
proposal distribution.  
To make this feasible, we fit a flexible proposal
$q_\theta$, e.g.\ a Gaussian mixture, to a set of ALDI samples
and subsequently perform IS with respect to $q_\theta$. Close to stationarity, the resulting mixture will satisfy the support condition~\eqref{eq:support-cond-F}, since
$\rho^*_{\delta,R}$ concentrates near the failure set $F$ for small
$(\delta,R)$ and assigns positive mass to all regions in $F$ where
$\rho_0$ is positive.  

Given $X_m\sim q_\theta$, the (unnormalised) IS weights are
\[
  w_m := \frac{\rho_0(X_m)}{q_\theta(X_m)},
\]
and the resulting self-normalised estimator is again
\[
  \widehat P_f
  = \frac{\sum_{m=1}^M\mathbf{1}_{\{G(X_m)\le 0\}}\,w_m}
         {\sum_{m=1}^M w_m}.
\]
This variant has the additional advantage that, once the proposal $q_\theta$ has
been fitted, an arbitrary number of independent samples can be generated at very
low cost. This enables the (self-normalized) IS estimator to achieve any prescribed accuracy
by drawing sufficiently many proposal samples. However, additional forward solves are thus needed to evaluate the IS estimator. 
Moreover, the separation between the ALDI phase (which generates the fitting
data) and the IS phase (which uses $q_\theta$) suggests that the overall work
can be balanced: ALDI does not need to be run until full convergence. As soon as
the ALDI cloud captures the relevant geometry of the failure region, the
subsequent IS step can be made as accurate as desired by increasing the number
of independent samples from $q_\theta$. This opens up the possibility of stopping
ALDI earlier while still reaching a target accuracy for $P_f$ at reduced
computational cost.

\begin{algorithm}[t]
\caption{Mixture-based IS fitted to ALDI samples}
\label{alg:aldi_mixture_is}
\begin{algorithmic}[1]
\Require prior $\rho_0$, limit-state $G$, smoothing parameters $\delta,R$,
         ALDI integrator, ensemble size $J$,
         fitting sample size $M_{\mathrm{fit}}$, IS sample size $M$,
         mixture size $K$
\Ensure estimate $\widehat P_f$ and effective sample size $\mathrm{ESS}$
\State run ALDI to stationarity and collect $M_{\mathrm{fit}}$ samples
       $\{X^{\mathrm{ALDI}}_m\}$
\State fit a parametric proposal $q_\theta$
       (e.g.\ a $K$-component Gaussian mixture) to the ALDI cloud
\State sample $X_1,\dots,X_M\sim q_\theta$
\For{$m=1,\dots,M$}
  \State $I_m:=\mathbf{1}_{\{G(X_m)\le 0\}}$
  \State $w_m:=\rho_0(X_m)/q_\theta(X_m)$
\EndFor
\State estimate
       $\displaystyle
        \widehat P_f=
        \frac{\sum_{m=1}^M I_m w_m}{\sum_{m=1}^M w_m}$
\State \Return $\widehat P_f$
\end{algorithmic}
\end{algorithm}



\subsection{Gradient-free ALDI method}
As an alternative formulation, a gradient-free version of the ALDI method can be used \cite{Nusken20}. One considers the sample cross-correlation matrix $\mathcal{D}(U)$ defined as
\begin{equation}
    \label{eq:crosscorrelation_matrix}
\mathcal{D}(X):= \frac{1}{J}\sum_{j=1}^J \left(x^{(j)}-m(X)\right) \left( \mathcal{G}(x^{(j)})-m(\mathcal{G}(X)) \right)^\top,
\end{equation}
where
\begin{equation}
    \label{eq:sample_mean2}
m(\mathcal{G}(X)):= \frac{1}{J}\sum_{j=1}^J \mathcal{G}\left(x^{(j)} \right),
\end{equation}
so that the \textit{gradient-free ALDI method} reads as
\begin{equation}
    \label{eq:ALDI_gradientfree}
\begin{split}
    dx^{(j)}_t & = -\left[ \mathcal{D}(X_t) R^{-1}\left( \mathcal{G}\left( x_t^{(j)} \right) -y_{obs}\right) +\mathcal{C}(X_t)P_0^{-1}\left(  x_t^{(j)}-m_0\right) \right]dt \\
    & \qquad + \frac{d+1}{J}\left(  x_t^{(j)}-m(X_t)\right)dt + \sqrt{2} \mathcal{C}^{1/2}(X_t)dW^{(j)}_t,
\end{split}
\end{equation}

\begin{remark}
    In our case, the forward map $\mathcal{G}=\tilde{G}$ and the observational noise $\eta$ are real-valued. Moreover, $y_{obs}=0$, which simplifies the computations. In fact,
\begin{equation}        \label{eq:ALDI_gradientfree_particular}
\begin{split}
dx_t^{(j)} & = -\left[ \frac{1}{R}\mathcal{D}(X_t)\tilde{G}\left(x_t^{(j)}\right) + \frac{1}{P_0}\mathcal{C}(X_t)\left( x_t^{(j)}-m_0\right) + \frac{d+1}{J}\left( x_t^{(j)}-m(X_t) \right)\right] dt \\
 & \qquad \qquad +\sqrt{2} \mathcal{C}^{1/2}(X_t)dW_t^{j}.
\end{split}
\end{equation}
\end{remark}

As shown in \cite{Nusken20}, the system \eqref{eq:ALDI_gradientfree} is also affine invariant, like its exact counterpart \eqref{eq:ALDI_gradientfree}. The drawback, however, is that a theoretical analysis as detailed in Theorem~\ref{THM:ALDI_properties} remains still an open problem.




\section{Illustrative Test Cases}
\label{SEC:numerics}

In this section, we assess the validity of the ALDI method for rare-event estimation by means of three illustrative examples. We first consider a classical example from structural reliability, introduced in
\cite{Katsuki94}. Secondly, we explore a simple linear ordinary differential equation with an unstable equilibrium of saddle type. Finally, we explore a dynamical model of interacting point-vortices.



\subsection{A convex limit-state function}
\label{SUBSEC:convex}
We consider the limit-state function
\begin{equation}
    \label{eq:convex_example}
    G(x_1,x_2)
    = 0.1 (x_1 - x_2)^2 - \frac{1}{\sqrt{2}}(x_1 + x_2) + 2.5,
\end{equation}
defines a convex failure domain.  As shown in \cite{Pappaioannou16}, the
distribution of $(x_1,x_2)$ conditioned on $\{G \le 0\}$ is unimodal {when the prior distribution is standard normal}, making this
a convenient test case for illustrating the behaviour of ALDI. In the following numerical experiments, we fix the smoothing parameter $\delta=0.001$, and the prior distribution $\mu_0\sim \mathcal{N}\left((0,0), Id \right)$. We solve ALDI numerically using the Euler-Maruyama method with a time step $\Delta \tau=0.001$ on an integration time window $\tau=10$. We used $8$ Gaussian components in the importance sampling step.


\paragraph{Numerical results}
We illustrate the behaviour of ALDI on the convex limit--state function
\eqref{eq:convex_example}. Figure~\ref{fig:convex_scatter} displays the final ALDI ensembles for three
ensemble sizes $J\in\{100,1000,10000\}$ (see panels (a), (b), and (c), respectively) and two values of the observational noise
parameter $R\in\{0.1,0.01\}$.
Smaller values of $R$ lead to a stronger concentration of particles near the
failure boundary $\{G=0\}$, while larger values of $R$ result in a more diffuse
distribution dominated by the prior.
Increasing the ensemble size $J$ substantially reduces sampling noise and
reveals the geometry of the target distribution more clearly.
Both gradient-based and gradient-free variants of ALDI exhibit the same
qualitative behaviour, with differences primarily visible for small ensembles. The reference solution is $P_f=4.21\times 10^{-3}$, cf.~\cite{Pappaioannou16} and verified by Monte Carlo.

\begin{figure}[h!]
    \centering
    \begin{overpic}[width=\linewidth]{./Convex_part1/convex__scatter.png}
           \put(15,-2){\footnotesize(a)}
           \put(48,-2){\footnotesize(b)}
           \put(81,-2){\footnotesize(c)}

           \put(15,33){\footnotesize$J=100$}
           \put(48,33){\footnotesize$J=1000$}
           \put(80,33){\footnotesize$J=10000$}
    \end{overpic}
    \caption{Final ALDI ensembles for the convex example.
    Columns correspond to ensemble sizes $J\in\{100,1000,10000\}$.
    Colours indicate gradient-based and gradient-free ALDI for
    $R=0.1$ and $R=0.01$.
    The red curve denotes the failure boundary $G=0$.}
    \label{fig:convex_scatter}
\end{figure}

The effect of the ensemble size $J$ is shown in
Figure~\ref{fig:size_and_failed} as a comparison between Alg.~\ref{alg:aldi_product} and Alg.~\ref{alg:aldi_mixture_is}.  
Although the identity $P_f = Z^*_{\delta,R}\,\mathbb P_{\rho^*_{\delta,R}}(X\in F)$
is exact and motivates Alg.~\ref{alg:aldi_product}, estimating the normalising constant
$Z^*_{\delta,R}$ directly from stationary ALDI samples via
\[
\widehat Z^*_{\delta,R}
=\left(\frac1M\sum_{m=1}^M \exp\!\Big(\tfrac{1}{2R}\tilde G_\delta(X_m)^2\Big)\right)^{-1}
\]
is typically numerically unstable. The
factor $\exp(\tilde G_\delta^2/(2R))$ can be extremely large away from the failure
set, so that rare excursions of the ALDI chain dominate the empirical average and
lead to very high variance. In addition, any deviation from stationarity introduces
bias, since the inverse-moment identity holds only under $\rho^*_{\delta,R}$. Indeed, if the empirical ALDI distribution differs from $\rho^*_{\delta,R}$,
the inverse-moment estimator converges to an expectation under the transient
distribution rather than $(Z_{\delta,R}^{*})^{-1}$. In addition, time discretization introduces a bias in the invariant measure of
the numerical ALDI scheme, which further degrades the accuracy of the normalisation
constant estimator.
For these reasons, we focus in the following on Alg.~\ref{alg:aldi_mixture_is}, which avoids estimating
$Z^*_{\delta,R}$ altogether by working with an explicit fitted proposal.
As expected, the estimation error decreases with increasing $J$.


\begin{figure}[h!]
    \centering
    \begin{subfigure}[t]{0.48\textwidth}
        \centering
        \begin{overpic}[width=\textwidth]{./Convex_part1/convex_Error_vs_ALDI_Ensemble_size_with_no_gaussian_fitting.png}
            \put(50,-2){\footnotesize $J$}
            \put(-5,30){\footnotesize \begin{turn}{90} $\lvert \widehat P_f- P_{\mathrm{ref}}\rvert$\end{turn}}
        \end{overpic}
        \caption{}
    \end{subfigure}
    \hfill
    \begin{subfigure}[t]{0.48\textwidth}
        \centering
        \begin{overpic}[width=\textwidth]{./Convex_part1/convex__Error_vs_ALDI_Ensemble_size.png}
         \put(50,-2){\footnotesize $J$}
            \put(-5,30){\footnotesize \begin{turn}{90} $\lvert \widehat P_f- P_{\mathrm{ref}}\rvert$\end{turn}}
        \end{overpic}
        \caption{}
    \end{subfigure}
    \caption{Dependence of the ALDI estimator on the ensemble size $J$. In (a), the results of Alg.~\ref{alg:aldi_product} are shown, in (b), the results of Alg.~\ref{alg:aldi_mixture_is} are shown for the corresponding number of ALDI samples $J$ in the IS estimator.}
    \label{fig:size_and_failed}
\end{figure}


Next, we investigate the influence of the smoothing parameter $\delta$.
For some values $\delta\in[10^{-4},10^{-1}]$, we compute the estimator $\widehat P_f$
using ALDI ensembles of size $J=100$ and $J=1000$. 
Figure~\ref{fig:delta_and_obs} (a) shows that the estimation error is essentially
flat for sufficiently small $\delta$, confirming the theoretical independence
of the limiting measure with respect to $\delta$.
For larger values of $\delta$, a mild deterioration in accuracy can be observed,
especially for smaller ensemble sizes.

Figure~\ref{fig:delta_and_obs} (b) illustrates the dependence on the observational
noise $R$.
Larger values of $R$ reduce the selectivity of the likelihood term
$\exp(-\tilde G^2/(2R))$, leading to moderately increased errors.
This effect is more pronounced for small ensembles and is mitigated as $J$
increases.

\begin{figure}[h!]
    \centering
    \begin{subfigure}[t]{0.48\textwidth}
        \centering
        \begin{overpic}[width=\textwidth]{./Convex_part1/convex__Error_vs_smoothness.png}
            \put(55,0){\footnotesize$\delta$}
            \put(-5,30){\footnotesize \begin{turn}{90} $\lvert \widehat P_f- P_{\mathrm{ref}}\rvert$\end{turn}}
        \end{overpic}
        \caption{}
    \end{subfigure}
    \hfill
    \begin{subfigure}[t]{0.48\textwidth}
        \centering
        \begin{overpic}[width=\textwidth]{./Convex_part1/convex__Error_vs_obs_cov.png}
            \put(55,0){\footnotesize$R$}
            \put(-5,30){\footnotesize \begin{turn}{90} $\lvert \widehat P_f- P_{\mathrm{ref}}\rvert$\end{turn}}
        \end{overpic}
        \caption{}
    \end{subfigure}
    \caption{Influence of the smoothing parameter $\delta$ (a) and the observational
    noise level $R$ (b) on the ALDI estimator for the convex example.
    The estimate is computed based on \ref{alg:aldi_mixture_is} with $10^3$ samples and $8$ mixture components fitted to each ALDI output.}
    \label{fig:delta_and_obs}
\end{figure}



The number of particles located inside the failure domain increases
with $J$, indicating improved exploration of the region of the rare event and reduced
weight degeneracy in the subsequent importance-sampling step, cp. Figure~\ref{fig:gm_components}.
We further examine the influence of the number of Gaussian mixture components
used in the importance-sampling proposal.
Figure~\ref{fig:gm_components} shows that increasing the number of components
generally improves accuracy for sufficiently large ensembles.
For very small ensembles, however, the estimator becomes more variable,
reflecting the increased difficulty of fitting expressive mixture models from
limited data.
\begin{figure}[h!]
    \centering
    \begin{subfigure}[t]{0.395\textwidth}
        \centering
         \centering
        \begin{overpic}[width=\textwidth]{./Convex_part1/convex__Number_of_failed_particles.png}
        \put(50,-2){\footnotesize$J$}
        \put(-2,25){\footnotesize \begin{turn}{90} $\#$ particles in $F$\end{turn}}
        \end{overpic}
        \caption{}
    \end{subfigure}
    \hfill
    \begin{subfigure}[t]{0.59\textwidth}
        \centering
        \begin{overpic}[width=\textwidth]{./Convex_part1/convex__Error_vs_num_of_gaussian_10e4_samples_.png}
        \put(20,0){$\#$ Gaussian mixture components}
        \put(2,18){\small \begin{turn}{90} $\lvert \widehat P_f- P_{\mathrm{ref}}\rvert$\end{turn}}
    \end{overpic}
        \caption{}
    \end{subfigure}
    \caption{In (a), the number of particles lying inside of the failure domain is shown. In (b), the absolute error {of the estimator $\hat{P}_f$} as a function of the number of Gaussian mixture
    components used in the importance-sampling proposal Alg.~\ref{alg:aldi_mixture_is} with $10^3$ samples.}
    \label{fig:gm_components}
\end{figure}

We finally assess the accuracy of the subsequent importance sampling step based
on the ALDI-generated proposal.
Figure~\ref{fig:is_error} reports the absolute error of the IS estimator as a
function of the total number of importance samples.
For all ensemble sizes, the IS error decays at the expected Monte--Carlo rate,
indicating that the fitted Gaussian mixture provides a sufficiently accurate
proposal distribution.
Larger ALDI ensembles lead to uniformly smaller IS errors, confirming that a
better approximation of the rare-event distribution yields more balanced
importance weights and improved sampling efficiency.
\begin{figure}[h!]
    \centering
    \begin{overpic}[width=0.7\textwidth]{./Convex_part1/convex_8_Error_vs_IS.png}
        \put(5,20){\footnotesize \begin{turn}{90}
        $\lvert \widehat P_f^{\mathrm{IS}}- P_{\mathrm{ref}}\rvert$
        \end{turn}}
        \put(40,8){\footnotesize $\#$ IS samples}
    \end{overpic}
    \caption{Absolute error of the IS estimator in Alg.~\ref{alg:aldi_mixture_is} as a function of
    the number of IS samples for different ALDI ensemble sizes.
    Larger ensembles yield more accurate proposals and faster error decay.}
    \label{fig:is_error}
\end{figure}

\paragraph{A rarer-event regime beyond Monte Carlo.}
We conclude with a more challenging variant of the convex example in which the
failure event is significantly rarer.
In this regime, crude Monte Carlo sampling becomes ineffective: within feasible
computational budgets, either no failure samples are observed or the resulting
estimators exhibit prohibitively large variance.
For this reason, Monte Carlo reference values are not reported in what follows. In the following, we fixed the prior distribution as Gaussian with mean $(-2,-2)$ and covariance matrix $0.8 Id$, and solved ALDI using the Euler-Maruyama method with time step $\Delta\tau=0.0005$ up to time $\tau=20$. The smoothing parameter was taken $\delta=0.001$ and $2$ Gaussian components in the importance sampling step.

Figure~\ref{fig:rare_scatter} shows the final ALDI ensembles for three ensemble
sizes and two noise levels.
Despite the increased rarity of the event, ALDI continues to concentrate
particles close to the failure boundary, thereby producing informative proposal
distributions for importance sampling.
Both gradient-based and gradient-free variants exhibit stable behaviour, with
larger ensembles leading to sharper resolution of the rare-event region.

\begin{figure}[h!]
    \centering
    \begin{overpic}[width=\linewidth]{./new_convex_m2_v0p8/convex__scatter.png}
           \put(15,-2){\footnotesize(a)}
           \put(48,-2){\footnotesize(b)}
           \put(81,-2){\footnotesize(c)}

           \put(15,33){\footnotesize$J=100$}
           \put(48,33){\footnotesize$J=500$}
           \put(81,33){\footnotesize$J=1000$}
    \end{overpic}
    \caption{Final ALDI ensembles in a rarer-event regime.
    Columns correspond to ensemble sizes $J=100,500,1000$.
    Colors indicate gradient-based and gradient-free ALDI for
    $R=0.01$.
    The red curve denotes the failure boundary $G=0$.}
    \label{fig:rare_scatter}
\end{figure}

We next assess the performance of the subsequent importance sampling step.
Figure~\ref{fig:rare_is} reports the estimated failure probability and the
corresponding variance as functions of the number of importance samples.
In contrast to standard Monte Carlo sampling, the ALDI-based importance sampling estimator remains stable in the rare-event regime. Increasing the ALDI ensemble size leads to a systematic reduction in the estimated variance, indicating that improved proposal quality directly translates into more efficient importance sampling. In contrast, sampling from the prior fails to capture the rare-event region altogether and therefore does not provide a meaningful estimator.

\begin{figure}[h!]
    \centering
    \begin{subfigure}[t]{0.48\textwidth}
        \centering
        \begin{overpic}[width=\textwidth]{./new_convex_m2_v0p8/Probability_vs_IS_samples_with_MC.png}
        \put(37,-3){\footnotesize $\#$ IS samples}
        \put(-7,20){\begin{turn}{90} \footnotesize$\hat{P}^{IS}_f$\end{turn}}
        \end{overpic}
        \caption{}
    \end{subfigure}
    \hfill
    \begin{subfigure}[t]{0.48\textwidth}
        \centering
        \begin{overpic}[width=\textwidth]{./new_convex_m2_v0p8/variance_vs_IS_samples_with_MC.png}
         \put(37,-3){\footnotesize $\#$ IS samples}
        \put(-5,20){\begin{turn}{90} \footnotesize$\mathrm{Var}(\hat{P}^{IS}_f)$\end{turn}}
        \end{overpic}
        \caption{}
    \end{subfigure}
    \caption{Importance sampling diagnostics in the rarer-event regime. The ALDI-generated proposal enables stable probability estimation (see (a)) and
    controlled variance decay (see (b)), even when crude Monte Carlo fails.}
    \label{fig:rare_is}
\end{figure}
\subsection{A hyperbolic saddle}
\label{SUBSEC:saddle}
We next consider a simple linear example that captures some of the dynamical
features underlying the point-vortex model discussed later in
Section~\ref{SUBSEC:vortex}.  Let $(x(t),y(t))$ solve
\begin{equation}
\label{eq:saddle_eq}
\begin{aligned}
    \dot{x} &= -\lambda x,\\
    \dot{y} &= \mu y,
\end{aligned}
\end{equation}
with $\lambda,\mu>0$.  The origin is a saddle equilibrium with one-dimensional
stable and unstable manifolds given by the horizontal and vertical axes,
respectively.  The explicit solution for initial data $(x_0,y_0)$ is
\begin{equation}
\label{eq:saddle_sols}
    x(t)=x_0 e^{-\lambda t}, 
    \qquad 
    y(t)=y_0 e^{\mu t}.
\end{equation}

Unless $y_0=0$, the trajectory escapes exponentially fast along the unstable
direction, given by the vertical axis.  We consider the rare event that a trajectory remains close to the
origin over a finite time window $[0,T]$.  This is quantified by the observable
\begin{equation}
    \label{eq:obs1_saddle}
    A(x_0,y_0)
    := \frac{1}{T} \int_0^T \bigl(x(s)^2 + y(s)^2 \bigr)\, ds.
\end{equation}
A short computation gives the closed-form expression
\[
A(x_0,y_0)
= \frac{1}{T}
\left[
  \frac{1-e^{-2\lambda T}}{2\lambda}\, x_0^2
  +
  \frac{e^{2\mu T}-1}{2\mu}\, y_0^2
\right],
\]
so that the rare-event set $\{A<r\}$, for $r\ll 1$, is the solid ellipse
\begin{equation}
\label{eq:failure_saddle_ellipse}
\mathcal{E}
=
\left\{
(x_0,y_0)\in\mathbb{R}^2:\ 
\left(\frac{1-e^{-2\lambda T}}{2\lambda r T}\right) x_0^2
+
\left(\frac{e^{2\mu T}-1}{2\mu r T}\right) y_0^2 
< 1
\right\}.
\end{equation}

We define the limit-state function
\begin{equation}
    \label{eq:LSF_saddle}
    G(x_0,y_0) = A(x_0,y_0) - r.
\end{equation}
For Gaussian prior samples $(x_0,y_0)\sim \mathcal{N}(m,\sigma^2 I)$, the
probability of failure is
\[
P_f
=
\iint_{\mathcal{E}} 
\rho_0(x,y)\,dx\,dy,
\qquad
\rho_0(x,y)
=
\frac{1}{2\pi\sigma^2}
\exp\!\left(
-\frac{(x-m_x)^2 + (y-m_y)^2}{2\sigma^2}
\right).
\]

This example is particularly convenient because the failure set is an ellipse
and $P_f$ can again be computed numerically to high accuracy.  This allows us to assess the performance of ALDI in a dynamical setting featuring unstable behaviour, where the rare event corresponds to trajectories remaining close to the saddle point.

The corresponding potential is
\begin{equation}
    \label{eq:potential_saddle}
    \Phi(x_0,y_0)
    =
    -\log \rho_0(x_0,y_0)
    +
    \frac{1}{2R}\,\tilde{G}(x_0,y_0)^2,
\end{equation}
where $\tilde G=\max\{0,G\}$.  
The potential reads
\begin{equation}\label{eq:Phi_saddle}
  \Phi(x_0,y_0)
  = -\log \rho_0(x_0,y_0)
    + \frac{1}{2R}\,\tilde G(x_0,y_0)^2,
\end{equation}
where $\tilde G=\max\{0,G\}$ and $G(x_0,y_0)=A(x_0,y_0)-r$ with
\begin{equation}\label{eq:A_saddle}
A(x_0,y_0)
= \frac{1}{T}\Bigg(
   \frac{1-e^{-2\lambda T}}{2\lambda}x_0^2
 + \frac{e^{2\mu T}-1}{2\mu}y_0^2
  \Bigg).
\end{equation}
The gradient of $G$ is computed explicitly as
\begin{equation}\label{eq:gradG_saddle}
\nabla G(x_0,y_0)
 = \frac{1-e^{-2\lambda T}}{\lambda T}\,x_0\,\mathbf e_x
 + \frac{e^{2\mu T}-1}{\mu T}\,y_0\,\mathbf e_y.
\end{equation}
Hence,
\begin{equation}\label{eq:gradPhi_saddle}
\nabla\Phi(x_0,y_0)
 = \frac{1}{\sigma^2}
   \begin{pmatrix}x_0-m_x\\ y_0-m_y\end{pmatrix}
 + \frac{\tilde G(x_0,y_0)}{R}\,
   \nabla G(x_0,y_0).
\end{equation}


\paragraph{Numerical results.}
We now illustrate the behaviour of ALDI for the saddle-type example introduced
above. In the following experiments, we fix $\lambda=\mu=1$, the time horizon $T=$ 1, the tolerance $r=0.5$, and the prior distribution to be normally distributed with mean $[-2,-2]$ and covariance matrix $0.5 Id$. Additionally, we solve ALDI numerically using the Euler-Maruyama method, with time step $\Delta \tau=\frac{0.001}{4}$ up to time $\tau=5$. In the importance sampling step, the number of Gaussian components used was 1.

Since the failure set $\{A<r\}$ is an ellipse and its probability under the
Gaussian prior can be computed accurately, this setting provides a controlled
benchmark for assessing the estimator $\widehat P_f$ in a dynamical situation
characterized by instability along the unstable manifold. All failure probability estimates in this section are computed using the mixture-based importance sampling estimator (Alg. \ref{alg:aldi_mixture_is}).

Figure~\ref{fig:saddle_scatter} shows representative ALDI ensembles for different ensemble
sizes $J\in\{100,500,1000\}$ and noise levels $R\in\{1,0.1,0.01\}$. As $R$ decreases,
the invariant measure concentrates increasingly near the failure ellipse, reflecting
the stronger selectivity of the likelihood. Increasing $J$ reduces sampling noise and
leads to a visibly higher density of particles inside the failure domain. While both
gradient-based and gradient-free ALDI capture the geometry of the rare-event set,
the gradient-based variant produces a sharper concentration near the boundary
$\{G=0\}$, particularly for smaller ensemble sizes.

\begin{figure}[h!]
    \centering
    \begin{overpic}[width=\textwidth]{new_saddle_m1_v0p5/hyp__scatter.png}
    \put(13,33){\small $J=100$}
    \put(45,33){\small$J=500$}
    \put(77,33){\small $J=1000$}
    \end{overpic}
    \caption{ALDI ensembles for the saddle example for different ensemble sizes
    $J$ and noise levels $R$. The red ellipse indicates the failure set
    $\{A<r\}$. Again, the red ellipse corresponds to the set $\{G=0\}$}
    \label{fig:saddle_scatter}
\end{figure}




The number of particles located inside the failure domain is shown in
Figure~\ref{fig:saddle_mix}(a). As expected, this number increases with $J$ and is larger for smaller values of $R$, reflecting the stronger concentration of the ALDI
invariant measure near the failure ellipse. The effect is particularly pronounced for
the gradient-based dynamics, which generate more informative proposals for the
subsequent IS step.

Finally, Figure~\ref{fig:saddle_mix}(b) illustrates the dependence of the estimator
accuracy on the number of Gaussian mixture components used in the proposal.
For sufficiently large ensembles, increasing the number of components improves the
approximation of the ALDI cloud and reduces the estimation error. For smaller
ensembles, however, overly expressive mixtures lead to increased variability,
reflecting the difficulty of fitting complex models from limited data.



\begin{figure}[h!]
    \centering
    \begin{subfigure}[t]{0.395\textwidth}
        \centering
         \centering
        \begin{overpic}[width=\textwidth]{./new_saddle_m1_v0p5/hyp__Number_of_failed_particles.png}
        \put(50,-2){\footnotesize$J$}
        \put(-2,20){\footnotesize \begin{turn}{90} $\#$ particles in $F$\end{turn}}
        \end{overpic}
        \caption{}
    \end{subfigure}
    \hfill
    \begin{subfigure}[t]{0.59\textwidth}
        \centering
        \begin{overpic}[width=\textwidth]{./new_saddle_m1_v0p5/hyp__Error_vs_num_of_gaussian_components.png}
        \put(20,0){\footnotesize$\#$ Gaussian mixture components}
        \put(3,18){\footnotesize \begin{turn}{90} $\lvert \widehat P_f- P_{\mathrm{ref}}\rvert$\end{turn}}
    \end{overpic}
        \caption{}
    \end{subfigure}
    \caption{In (a), the number of particles lying inside of the failure domain is shown. In (b), the absolute error {of the estimator $\hat{P}_f$} as a function of the number of Gaussian mixture
    components used in the importance-sampling proposal Alg.~\ref{alg:aldi_mixture_is} with $10^3$ samples.}
    \label{fig:saddle_mix}
\end{figure}

We next examine the accuracy of the importance-sampling estimator obtained from
the ALDI-generated proposals. Figure~\ref{fig:saddle_is_error} reports the absolute
error of $\widehat P_f$ as a function of the number of IS samples. For all ensemble
sizes, the error decays at the expected Monte--Carlo rate, confirming that the fitted
Gaussian mixtures provide effective proposal distributions. 

\begin{figure}[h!]
    \centering
    \begin{overpic}[width=0.85\textwidth]{./new_saddle_m1_v0p5/saddle_1_Error_vs_IS.png}
        \put(5,15){\small \begin{turn}{90}
        $\lvert \widehat P_f^{\mathrm{IS}}- P_{\mathrm{ref}}\rvert$
        \end{turn}}
        \put(40,00){\small $\#$ IS samples}
    \end{overpic}
    \caption{Absolute error of the IS estimator in Alg.~\ref{alg:aldi_mixture_is} as a function of
    the number of IS samples for different ALDI ensemble sizes.}
    \label{fig:saddle_is_error}
\end{figure}


\paragraph{Summary.}
Across both test cases, the numerical results are in good agreement with the
theoretical predictions. ALDI yields accurate estimates of the failure
probability, with errors decreasing as the ensemble size $J$ increases. Once
the smoothing parameter $\delta$ is sufficiently small, the estimator is largely
insensitive to its precise value, whereas smaller values of $R$ sharpen the
auxiliary distribution around the rare-event region and improve localization of
the failure set. Consistently across both examples, the number of particles
lying inside the failure domain increases with $J$ and is higher for smaller
$R$, indicating a better-adapted proposal for subsequent importance sampling.

Regarding the two estimation strategies, Alg.~\ref{alg:aldi_product} provides a useful
conceptual link between the biased ALDI measure and the original failure
probability through a product representation. However, the numerical
experiments confirm that directly estimating the associated normalisation
constant is unstable in practice and sensitive to incomplete stationarity of
the ALDI dynamics. In contrast, Alg.~\ref{alg:aldi_mixture_is} avoids this difficulty by fitting
an explicit proposal distribution to the ALDI ensemble and applying standard
importance sampling. This approach yields robust and systematically improving
estimates as $J$ increases, and its performance is directly correlated with the
quality of the ALDI-generated samples.

Overall, the two examples demonstrate that ALDI provides a reliable mechanism
for constructing effective proposals for rare-event estimation across problems
of different geometric and dynamical structure, while Alg.~\ref{alg:aldi_mixture_is} offers a
practically stable and accurate estimator in regimes where direct normalisation
is infeasible.



\subsection{Atmospheric blockings and rare events}
\label{SUBSEC:vortex}

An \textit{atmospheric blocking} is a quasistationary structure in the atmospheric flow over a large region. Their timescale ranges from several days to weeks, and can have a drastic impact in the local weather \cite{Kuhlbrodt00}. It becomes then of crucial relevance to estimate the probability of duration of an atmospheric blocking to remain over a long time window. Events of this nature are often captured by means of a real-valued function $G$  \cite{Ragone18,Ragone20}.

\subsubsection{Point-vortex dynamics}
A simple description of atmospheric flows at synoptic and planetary scales is given by the standard planar Euler equations, where the flow is assumed to be inviscid and incompressible \cite{Muller15}. Its expression in vorticity form is thus
\begin{equation}
\label{eq:Euler_vorticity}  
    \frac{\partial \bm{\omega}}{\partial t}+\bm{u}\cdot\nabla\bm{\omega}=\bm{\omega}\cdot \nabla \bm{u},
\end{equation}
where $\bm{u}(t,x)$ and $\bm{\omega}(t,x)$ are the velocity and vorticity fields of the flow on $x\in\mathbb{R}^2$ at time $t\geq0$, respectively.

Furthermore, assuming the vorticity is concentrated in $N$ disjoint points $X_j\in\mathbb{R}^2$, i.e.
\[
    \bm{\omega}(t,x)=\sum_{j=1}^N \Gamma_j \delta_{X_j(t)}(x),
\]
where $\delta_X$ denotes the Dirac delta distribution supported on $X$, one can derive an ODE system for the position of the point-vortices $X_j=(x_j,y_j)$, which reads \cite{Newton01}
 \begin{equation}
    \label{eq:Nvortices_plane}
        \begin{split}
            \dot{x}_j & =-\frac{1}{2\pi} \sum_{i\neq j} \frac{\Gamma_i(y_j-y_i)}{l_{ij}^2} \\
            \dot{y}_j & = \frac{1}{2\pi} \sum_{i\neq j} \frac{\Gamma_i(x_j-x_i)}{l_{ij}^2},
        \end{split}
    \end{equation}
    where $l_{ij}=\Vert X_i-X_j \Vert$ are the \textit{intervortical distances}. The system \eqref{eq:Nvortices_plane} is referred to as the \textit{$N$-vortex problem}. 
    
\begin{remark}
The derivation of the point-vortex dynamics turns out to be a good approximation of \eqref{eq:Euler_vorticity}, whenever the initial vorticity $\bm{\omega}(0,x)$ is supported on $N$ disjoint regions $\Lambda_j^{\varepsilon}$, contained in circles of radius $\varepsilon$ centred on $X_j(0)$, for $j=1,\ldots, N$. Indeed, under suitable assumptions, it has been shown that the corresponding solutions of \eqref{eq:Euler_vorticity}, converge weakly as $\varepsilon\rightarrow0$ to those given by \eqref{eq:Nvortices_plane}. We refer to \cite[Theorem 2.1]{Marchioro93} for the precise formulation of the statement. More recently, a connection of a stochastically perturbed version of \eqref{eq:Nvortices_plane} and the deterministic Navier-Stokes equation has been established, in the large ensemble limit $N\rightarrow\infty$ \cite{Flandoli24}
\end{remark}

The $N$-vortex problem has a noncanonical Hamiltonian structure. Indeed, if $\bm{X}=(X_1,\ldots,X_N)$ the system  \eqref{eq:Nvortices_plane} is written as
\[
\dot{\bm{X}} = \mathbb{J} \nabla \mathcal{H} (\bm{X}), \qquad \mathbb{J} = \left[ \begin{array}{cc}
    0 & \textup{diag}(\Gamma_1^{-1},\ldots,\Gamma_N^{-1}) \\
    -\textup{diag}(\Gamma_1^{-1},\ldots,\Gamma_N^{-1}) & 0
\end{array} \right]
\]
where the energy function $\mathcal{H}$ is given by
\begin{equation}
    \label{eq:Energy_Nvortex}
    \mathcal{H}(\bm{X}) = -\frac{1}{4\pi}\sum_{
          i\neq j} \Gamma_i\Gamma_j \ln l_{ij},
\end{equation}
and the sum is taken over all pairs $i,j\in\{1,\ldots,N\}$. For a comprehensive exposition of the dynamics of \eqref{eq:Nvortices_plane}, we refer to \cite{Newton01}.

\subsubsection{Atmospheric blockings as relative equilibria of point-vortices}

In \cite{Kuhlbrodt00}, blocking events are described by the dynamics of low-order point-vortex systems, typically when $N\leq4$. The mechanism is as follows: whenever a vortex configuration is relatively stable, any passive vortex $X_{N+1}$ (i.e. a vortex with circulation $\Gamma_{N+1}=0$) remains secluded in a large area surrounding the position of the vortices $X_1,\ldots, X_N$ for large time windows. An escape from such region can happen due to an external forcing or noisy fluctuations, see Figure~\ref{fig:passive_vortices}.  When the blocking is constituted by two (quasi)stationary vortices, we say that a \textit{high-over-low blocking} is present, see panel (a) in Figure~\ref{fig:passive_vortices}. In the case that it consists of a three point-vortex configuration, we call it an \textit{$\Omega$ blocking}, see panel (b) in Figure~\ref{fig:passive_vortices}.

\begin{figure}
     \centering
     \begin{subfigure}[b]{0.45\textwidth}
         \centering
          \begin{overpic}[width=\linewidth]{contour2.png}
           \put(50,-2){$x$}
      \put(3,45){$y$}
      \end{overpic}
      \caption{}
         \label{subfig:dipole_nonoise}
     \end{subfigure}
     \hfill
     \begin{subfigure}[b]{0.45\textwidth}
         \centering
          \begin{overpic}[width=\linewidth]{contour3.png}
           \put(50,-2){$x$}
      \put(3,45){$y$}
      \end{overpic}
      \caption{}
         \label{subfig:tripole_nonoise}
     \end{subfigure}
     \hfill
       \caption[Advected $N+1$-vortex problem]{Depiction of dipole and tripole configurations of point vortices in equilibrium. On panels (a) and (b), the energy landscapes for a passive vortex under the influence of two and three stationary point-vortices located at $X_j\in \mathbb R^2$ (marked as red crossed), respectively, are portrayed in solid lines. The Hamiltonian for the passive vortex $X=(x,y)$ under the influence of $N$ stationary vortices is given by $\mathcal{H}_{N+1}(x,y):= -\frac{1}{8\pi}\sum_{j=1}^N\ln\; \Vert X-X_j\Vert^2$, whose dynamics is perturbed with additive Brownian noise. Exemplary trajectories of such passive vortices, under the influence of additive Brownian noise, confirm the quasistationary region induced by the vortex configuration.}
    \label{fig:passive_vortices}
\end{figure}

In the case of the 2-vortex problem, the dynamics is well understood - in the case that the total circulation $\Gamma=\Gamma_1+\Gamma_2$ is nonzero, both vortices $X_1,X_2$ rotate around the \textit{centre of vorticity} $X_c=(\Gamma_1/\Gamma) X_1 + (\Gamma_2/\Gamma) X_2$ with constant angular velocity $\Gamma/(2\pi\Vert X_1(0)-X_2(0)\Vert^2)$, see for instance \cite[p. 39]{Newton01}. In the case $\Gamma=0$, the vortices move in a parallel fashion, with speed $\sqrt{\Gamma}/(2\sqrt{2}\pi \Vert X_1(0)-X_2(0)\Vert)$. In any case, all possible dynamics are \textit{relative equilibria}, in the sense that the intervortical distance $l_{12}$ remains constant for all $t>0$.

\textit{Fixed equilibria} in the three-vortex problem, i.e. configurations in which $\dot{X}_i=0$ for $i=1,2,3$, are characterised by the positions being collinear and that the circulations satisfy $\Gamma_1\Gamma_2 + \Gamma_2\Gamma_3+ \Gamma_3\Gamma_1=0$. A configuration is a relative equilibrium, if and only if it is an equilateral triangle \cite[Theorem 2.2.1]{Newton01}. In the case $\Gamma:=\Gamma_1+\Gamma_2+\Gamma_3\neq 0$, the configuration rotates around the centre of vorticity with angular velocity $\Gamma/2\pi l^2$, where $l=l_{ij}$. When $\Gamma=0$, the configuration translates in parallel with velocity $(\Gamma_1^2+\Gamma_2^2+\Gamma_3^2)^{1/2}/(2\sqrt{2}\pi l)$.

From the observations above, it is clear almost none of the possible configurations in low-order point vortex dynamics admits fixed equilibria, and the pictures in Figure~\ref{fig:passive_vortices} may be misleading as a representation of the real phenomenon. However, the angular velocity depends on the total circulation $\Gamma$, in such a way that it vanishes as $\Gamma\rightarrow0$ or as the intervortical distances $l\rightarrow\infty$. Therefore, a quasi-stationary regime is expected at large scales.

\subsubsection{Stability of relative equilibria and rare events}
By definition, the space of relative equilibria $\mathcal{R}$ is an invariant manifold for the dynamics induced by \eqref{eq:Nvortices_plane}. In mathematical terms, the manifold $\mathcal{R}$ reads
\begin{equation}
    \label{eq:relative_equilibria_manifold}
\mathcal{R}= \left\{ \bm{X}\in\mathbb{R}^{2N} \; : \;  \frac{d l_{ij}}{dt}=0 \textup{ for } i,j=1,\ldots,N \right\}.
\end{equation}

As mentioned above, any solution of the 2-vortex problem is a relative equilibrium, and therefore $\mathcal{R}$ is the whole phase space $\mathbb{R}^{4}$. However, this is not the case for the 3-vortex problem. For instance, the set of equilateral triangular configurations $\mathcal{R}_\Delta$ is an invariant submanifold of $\mathcal{R}$ with $\textup{dim}\mathcal{R}_\Delta=4$. Indeed, the first vortices $X_1$ and $X_2$ are arbitrary points in $\mathbb{R}^2$, while the third one has only two possibilities, corresponding to the two possible orientations of the triangle $\Delta X_1 X_2 X_3$. Therefore, $\mathcal{R}_\Delta$ consists of two mutually disjoint branches.


The linear stability of $\mathcal{R}$ is given by the eigenvalues of its linearisation, where two of them are $0$, and two more are purely imaginary (possibly zero as well). Hence its stability is determined by the remaining two, which implythe following result.
\begin{theorem}
The invariant manifold $\mathcal{R}$ is unstable (of saddle type) if and only if \begin{equation}
    \gamma:=\Gamma_1\Gamma_2 + \Gamma_2\Gamma_3 + \Gamma_3\Gamma_1 <0. 
\end{equation}
In such a case, $\mathcal{R}_\Delta$ is a \textit{normally hyperbolic invariant manifold} of \eqref{eq:Nvortices_plane} with one-dimensional stable and unstable directions. More specifically, let $\varphi_t$ be the dynamical system induced by the 3-vortex problem. Then, there is a splitting of the tangent bundle
\[ T\mathbb{R}^6\vert_{\mathcal{R}_\Delta} = \mathcal{N}^s \oplus T\mathcal{R}_\Delta \oplus\mathcal{N}^u, \]
and local manifolds $W^{s,u}_{loc}(M)$ such that 

\smallskip

\begin{enumerate}
    \item $W^{s,u}_{loc}(\mathcal{R}_\Delta)$ is fibred as
    \[W^{s,u}_{loc}(\mathcal{R}_\Delta)= \bigcup_{p\in \mathcal{R}_\Delta} \mathcal{F}^{s,u}(p),\]
    where $\mathcal{F}^{s,u}(p)$ is a one-dimensional manifold based on $p\in\mathcal{R}_\Delta$,
    \item $\mathcal{F}^{s,u}(p)$ is tangent to $\mathcal{N}^{s,u}$ at $p$,
    \item The fibres $\mathcal{F}^{s}(p)$ and $\mathcal{F}^{u}(p)$ are positively and negatively invariant, respectively, that is for every $p\in\mathcal{R}_\Delta$ and $t\geq 0$,
    \begin{align*}
        \varphi_{t}\left( \mathcal{F}^s(p) \right) &\subset \mathcal{F}^s\left( \varphi_t(p)\right), \\
         \varphi_{-t}\left( \mathcal{F}^u(p) \right) &\subset \mathcal{F}^s\left( \varphi_{-t}(p)\right),
    \end{align*}
    \item There exist $C_{s,u}, \lambda_{s,u}>0$ such that if $q_{s,u}\in\mathcal{F}^{s,u}$, then
    \begin{align*}
        \Vert \phi_t(p)-\phi_t(q_s) \Vert < C_se^{-\lambda_st} \\
        \Vert \phi_{-t}(p)-\phi_{-t}(q_u) \Vert < C_ue^{-\lambda_ut}
    \end{align*}
\end{enumerate}
\end{theorem}
\begin{proof}
    We refer to \cite{Aref09} for the full analysis of the linear stability of relative equilibria. The structure of the tangent bundle into its stable and unstable parts is given by standard invariant manifold theory of dynamical systems, see for instance \cite{Wiggins94}.
\end{proof}
As observed in \cite{Aref09}, we have the following result.
\begin{corollary}
\label{CORO:equilibria_zeroGamma}
If $\Gamma=0$, any equilateral triangular configuration is an unstable relative equilibrium.
\end{corollary}
\begin{proof}
    It follows directly from the relationship $\Gamma^2 = (\Gamma_1^2+\Gamma_2^2+\Gamma_3^2) + 2\gamma$.
\end{proof}

In \cite{Muller15,Hirt18}, $\Omega$ blockings have been described by relative equilibria of three point vortices, whose circulations are given by $\Gamma_1=\Gamma_2$ and $\Gamma_3=-2\Gamma_1$, so that the total circulation $\Gamma$ vanishes. As stated in Corollary~\ref{CORO:equilibria_zeroGamma}, these configurations are unstable, of saddle type, so that small perturbations from such constellations deviate exponentially fast, unless the perturbed initial condition lies on $W^s_{loc}(\mathcal{R}_\Delta)$. Since $\textup{dim}\;W^s_{loc}(\mathcal{R}_\Delta)= 5$, it is a Lebesgue null set, and therefore this situation does not happen in practice. The association of atmospheric blockings to invariant and dynamically unstable objects has been also reported in \cite{Faranda16}.

Instead, one can explore initial conditions which remain close to a relative equilibrium configuration for large time windows \cite{Ragone18,Ragone20}. In this case, given an energy value $H$, we consider the observable
\begin{equation}
    \label{eq:observable_releq}
    A(X_1,X_2,X_3):= \left\vert \cos\theta_1 -1/2\right\vert + \left\vert \cos\theta_2 -1/2\right\vert  + \left\vert\frac{ l_{12} + l_{23} + l_{31} }{3} -l(H)\right\vert,
\end{equation}
where $\theta_1$ is the angle between the vectors $X_2-X_1$ and $X_3-X_1$, $\theta_2$ is the angle betwwen the vectors $X_1-X_2$ and $X_3-X_2$, and $l(H)$ is the length of the unique equilateral configuration associated to the energy level $\mathcal{H}(X_1,X_2,X_3)=H$, that is $l(H)=e^{-\frac{4\pi H}{\gamma}}$. Clearly $A(\bm{X})\equiv A(X_1,X_2,X_3)=0$ if and only if $(X_1,X_2,X_3)\in\mathcal{R}_\Delta(H)$, i.e. if the points $X_1,X_2,X_3$ form an equilateral triangle of length $l(H)$. 

Additionally, we consider additive stochastic perturbations, as given by the SDE
\begin{equation}
    \label{eq:SDE_vortex}
d\bm{X}_t = \mathbb{J} \nabla\mathcal{H}(\bm{X}_t) + \sigma dW_t,
\end{equation}
with given initial condition $\bm{X}_0$. 
Here, $\sigma>0$ is the noise strength, and $W_t$ is a standard brownian motion in $\mathbb{R}^6$. We denote the stochastic process induced by \eqref{eq:SDE_vortex} as $X_t(\omega)$. For a given time window $T>0$, let
\begin{equation}
    \label{eq:averaged_observable}
\bar{A}_T(\bm{X},\omega):= \frac{1}{T}\int_0^TA\left( \bm{X}_s(\omega) \right)ds.
\end{equation}
For $r>0$ sufficiently small, the rare event consists of the set
\[
    \left\{ (\bm{X}_0,\omega) : A_T(\bm{X}_s(\omega))<r \right\}.
\]
Therefore, our proposed limit-state function $G$ is defined as
\begin{equation}
    \label{eq:limit-state function_vortex}
G_r(\bm{X}):= \bar{A}_T(\bm{X})-r.
\end{equation}


\paragraph{Numerical results}

In the upcoming numerical explorations, we fix the circulations $\Gamma_1=\Gamma_2=1$ and $\Gamma_3=-2$. We fix the initial condition at an energy level $\mathcal{H}(\mathbf{X})=1$ on an equilateral triangle configuration of length $l(1)$. The threshold value used in this case is, $r=0.25$. We solve the $3$-vortex problem using the Euler-Maruyama method, with step size $\Delta T=0.02$ on a time horizon $T=0.5$. Additionally, the ALDI method is solved numerically using a step size $\Delta \tau=0.0005$ on an integration timescale $\tau=2.5$. In the importance sampling step, the number of Gaussian components used was 1.

Rare events are defined in terms of anomalous vortex configurations,
characterized by unusually persistent intervortical distances over a prescribed
time window. Figure~\ref{fig:vortex_structure} shows representative vortex
trajectories at different times, highlighting the emergence of dynamically
coherent structures associated with rare-event behaviour. While typical
configurations rapidly separate, rare-event trajectories remain clustered and
exhibit a markedly different geometric evolution. Although crude Monte Carlo sampling is prohibitively expensive in this regime,
a high-accuracy reference value $P_{\mathrm{ref}}$ is available from a
computationally intensive benchmark calculation and is used solely for
validation purposes.

\begin{figure}[h!]
    \centering
    \includegraphics[width=\textwidth]{new_vortex/vortex_structure_by_time.png}
    \caption{Evolution of representative three-vortex configurations at
    different times when the total circulation $\Gamma=0$. The system starts on an equilateral configuration, which is unstable in the absence of noise. Rare-event trajectories remain dynamically coherent,
    whereas typical trajectories exhibit rapid separation from the equilateral configuration.}
    \label{fig:vortex_structure}
\end{figure}

This behaviour is further quantified in
Figure~\ref{fig:vortex_distance_time}, which displays the evolution of
intervortical distances for different time horizons. For longer horizons,
typical trajectories show strong separation, whereas rare-event trajectories
remain close over extended periods, confirming the dynamical relevance of the
chosen failure event.

\begin{figure}[h!]
    \centering
        \begin{overpic}[width=0.65\textwidth]{new_vortex/A_diagnostic_vs_time_0p5.png}
        \put(50,0){\footnotesize$t$}
        \put(-5,30){\footnotesize$A(\mathbf{X}_t)$}
        \end{overpic}    
        
        \caption{The observable $A$ evaluated along the trajectories $X_t$ sampled using ALDI, as a function of time for different time. Here $J=1000$. Rare-event trajectories have relatively low value of $A$, compared to typical behaviour. The rare event sampled by ALDI corresponds to the trajectories portrayed in orange and red, corresponding to those for which $G\leq -0.22$ (647 particles) and $G\leq 0$ (18 particles), respectively, showing that the event is no longer rare for the ALDI samples.} 
    \label{fig:vortex_distance_time}
\end{figure}

Figure~\ref{fig:vortex_mix} summarizes the behaviour of the ALDI-generated
proposals and the subsequent mixture-based importance sampling step for the
point-vortex problem. Panel~(a) shows the number of particles lying inside the
failure domain as a function of the ALDI ensemble size $J$. As in the previous
examples, this number increases monotonically with $J$, indicating that larger
ensembles provide a more faithful exploration of the dynamically relevant rare
event region. Despite the strong sensitivity of the vortex dynamics to initial
conditions, ALDI consistently generates a non-negligible fraction of samples in
the failure set, which is essential for constructing effective IS proposals.

Panel~(b) of Figure~\ref{fig:vortex_mix} displays the absolute error of the
mixture-based estimator $\widehat P_f$ as a function of the number of Gaussian
mixture components used in the proposal. For sufficiently large ensembles, the
estimation error decreases as the mixture becomes more expressive, reflecting an
improved approximation of the ALDI sample cloud. For smaller ensembles, however,
the error exhibits increased variability when too many components are used,
highlighting the trade-off between model flexibility and the amount of available
fitting data.


\begin{figure}[h!]
    \centering
    \begin{subfigure}[t]{0.395\textwidth}
        \centering
         \centering
        \begin{overpic}[width=\textwidth]{./new_vortex/vortex__Number_of_failed_particles.png}
        \put(50,-2){\footnotesize$J$}
        \put(-2,20){\footnotesize \begin{turn}{90} $\#$ particles in $F$\end{turn}}
        \end{overpic}
        \caption{}
    \end{subfigure}
    \hfill
    \begin{subfigure}[t]{0.59\textwidth}
        \centering
        \begin{overpic}[width=\textwidth]{./new_vortex/hyp__Error_vs_num_of_gaussian}
        \put(20,0){$\#$ Gaussian mixture components}
        \put(2,18){\small \begin{turn}{90} $\lvert \widehat P_f- P_{\mathrm{ref}}\rvert$\end{turn}}
    \end{overpic}
        \caption{}
    \end{subfigure}
    \caption{In (a), the number of particles lying inside of the failure domain is shown. In (b), the absolute error {of the estimator $\hat{P}_f$} as a function of the number of Gaussian mixture
    components used in the importance-sampling proposal Alg.~\ref{alg:aldi_mixture_is} with $10^3$ samples.}
    \label{fig:vortex_mix}
\end{figure}

The performance of the importance sampling correction is examined in
Figure~\ref{fig:vortex_is_error}, which reports the absolute error of the IS
estimator as a function of the number of importance samples. For all ensemble
sizes, the error decreases at the expected Monte--Carlo rate, confirming that the
mixture proposals fitted to the ALDI output yield stable and consistent IS
estimates. Larger ALDI ensembles lead to uniformly smaller errors, demonstrating
that improved localization of the rare-event region during the ALDI phase
directly translates into higher sampling efficiency in the IS step.
\begin{figure}[h!]
    \centering
    \begin{overpic}[width=0.85\textwidth]{./new_vortex/vortex__Error_vs_IS_size.png}
        \put(5,15){\small \begin{turn}{90}
        $\lvert \widehat P_f^{\mathrm{IS}}- P_{\mathrm{ref}}\rvert$
        \end{turn}}
        \put(40,00){\small $\#$ IS samples}
    \end{overpic}
    \caption{Absolute error of the IS estimator in Alg.~\ref{alg:aldi_mixture_is} as a function of
    the number of IS samples for different ALDI ensemble sizes.}
    \label{fig:vortex_is_error}
\end{figure}

\section{Conclusions}
\label{SEC:conclusions}
We have proposed an ALDI approach for the
efficient sampling of rare events. Our analysis shows that, after introducing a smooth
approximation of the limit-state function, the resulting posterior distribution
retains the correct small-noise limit and is amenable to sampling by ALDI. The
method naturally adapts to anisotropies
in the posterior induced by low-probability geometries.

The numerical experiments demonstrate that ALDI captures the qualitative and
quantitative features of rare-event sets across a hierarchy of models: from
convex and saddle-type finite-dimensional examples, and finally to point-vortex models for atmospheric
blockings. In all cases, the sampler concentrates around the relevant
near-critical manifolds and yields accurate proposal distributions for
self-normalised importance sampling.

Overall, ALDI provides an effective and robust tool for rare-event estimation in
complex dynamical systems, particularly in settings where gradients are
unreliable or unavailable and where the geometry of the event is highly
anisotropic. Future work includes quantitative convergence analysis of ALDI,
extensions to fully time-dependent PDEs and stochastic systems,
multi-level or adaptive variants that reduce the cost of sampling,
and combinations with likelihood-free or surrogate-based approaches
for high-dimensional geophysical applications.

%\appendix

\section*{Code availability}
The implementation of the methods and all scripts used to produce the results in this paper are publicly available at \cite{Chakraborty25}.

\section*{Acknowledgments}
The work has been funded by Deutsche Forschungsgemeinschaft (DFG) through
the grant CRC 1114 ‘Scaling Cascades in Complex Systems’ (project number 235221301,
project A02 - Multiscale data and asymptotic model assimilation for atmospheric flows).

The authors would like to thank Antonia Liebenberg for providing parts of the code used in the numerical experiments. The authors also thank Robert Gruhlke for insightful discussions.





\bibliography{references}
\bibliographystyle{siamplain}


\end{document}
